\section{Модель и радиальное сведение}
\label{compact:assumptions}\label{sec:model}
Пусть $Q\subset\mathbb R^d$ есть нетривиальный выпуклый компакт,
$f\in C^2$ в окрестности $Q$ и $x^*\in\arg\min_Q f$.
Обозначим
\begin{equation}\label{eq:gap}
 \delta(x)=f(x)-f(x^*),\qquad P(x)=\|\nabla f(x)\|,
 \qquad G(x)=\max_{v\in Q}\langle\nabla f(x),x-v\rangle.
\end{equation}
Полный зазор $G$ использует всё множество $Q$.
В разделах о сертификатах предполагаются известными
$m>0$, $H_0,H_1\ge0$, для которых
\begin{equation}\label{boundarygap:assumptions}
 \begin{aligned}
 \langle\nabla f(z)-\nabla f(x^*),z-x^*\rangle
   &\ge m\|z-x^*\|^2,\\
 \|\nabla^2 f(z)\|&\le H_0+H_1\delta(z),\qquad z\in Q.
 \end{aligned}
\end{equation}
Первое условие называется сдвинутым RSI. Выпуклость $f$ и
стационарность $x^*$ не требуются. Для результатов через норму
достаточно, чтобы $Q$ содержало каждый отрезок $[x^*,x]$, а
$f$ была $C^2$ в окрестности этих отрезков.

При $x\ne x^*$ положим $r=\|x-x^*\|$, $e=(x-x^*)/r$,
$u(s)=f(x^*+se)-f(x^*)$ на $[0,r]$. Условная минимальность и
\eqref{boundarygap:assumptions} дают
\begin{equation}\label{compact:radial-data}
 q:=u'(0)\ge0,\qquad u'(s)\ge q+ms,
 \qquad |u''(s)|\le H_0+H_1u(s).
\end{equation}
Дифференцируя неравенство для $u'$ в нуле, получаем $H_0\ge m$.
Для $p=u'(r)$ имеем
\begin{equation}\label{compact:observed-caps}
 p\le P(x),\qquad rp\le G(x),\qquad
 r\le\min\{P(x)/m,\sqrt{G(x)/m}\}.
\end{equation}
Поэтому нулевой зазор или нулевая норма влекут $x=x^*$.
Норма в условном минимуме при этом может быть положительной.

Утверждения о точности ниже относятся к супремумам по функциям
и областям с заданными константами. Область может меняться.
Точность для одной фиксированной области не утверждается.
Гладкие последовательности могут приближать супремум, не достигая его.
